% ============================================
% Documento de Ejemplo - LaTeX
% Compilar con: pdflatex ejemplo.tex
% ============================================
\documentclass[12pt, a4paper]{article}

% ── Paquetes ──────────────────────────────────
\usepackage[utf8]{inputenc}        % Soporte UTF-8 (tildes, ñ, etc.)
\usepackage[spanish]{babel}        % Idioma español
\usepackage[T1]{fontenc}           % Codificación de fuentes
\usepackage{lmodern}               % Fuente Latin Modern (mejor calidad)
\usepackage{amsmath, amssymb}      % Matemáticas avanzadas
\usepackage{graphicx}              % Imágenes
\usepackage{xcolor}                % Colores
\usepackage{hyperref}              % Enlaces clickeables
\usepackage{geometry}              % Márgenes
\usepackage{fancyhdr}              % Encabezados/pies de página
\usepackage{booktabs}              % Tablas profesionales
\usepackage{enumitem}              % Listas personalizadas
\usepackage{listings}              % Código fuente

% ── Configuración de página ──────────────────
\geometry{
    top=2.5cm,
    bottom=2.5cm,
    left=3cm,
    right=3cm,
    headheight=14.5pt
}

% ── Configuración de hyperref ────────────────
\hypersetup{
    colorlinks=true,
    linkcolor=blue!70!black,
    urlcolor=blue!70!black,
    citecolor=green!50!black,
    pdftitle={Documento de Ejemplo},
    pdfauthor={Tu Nombre}
}

% ── Encabezados y pies de página ─────────────
\pagestyle{fancy}
\fancyhf{}
\fancyhead[L]{\textit{Documento de Ejemplo}}
\fancyhead[R]{\textit{\today}}
\fancyfoot[C]{\thepage}
\renewcommand{\headrulewidth}{0.4pt}

% ── Configuración de código fuente ───────────
\lstset{
    basicstyle=\ttfamily\small,
    keywordstyle=\color{blue!70!black}\bfseries,
    commentstyle=\color{green!50!black}\itshape,
    stringstyle=\color{red!70!black},
    numbers=left,
    numberstyle=\tiny\color{gray},
    frame=single,
    breaklines=true,
    captionpos=b
}

% ── Datos del documento ──────────────────────
\title{
    \textbf{Documento de Ejemplo en \LaTeX}\\[0.5em]
    \large Una guía rápida de lo que puedes hacer
}
\author{Tu Nombre Aquí}
\date{\today}

% ============================================
\begin{document}

\maketitle
\tableofcontents
\newpage

% ── SECCIÓN 1 ──────────────────────────────
\section{Introducción}

¡Bienvenido a tu primer documento \LaTeX! Este archivo de ejemplo te muestra las principales funcionalidades que puedes usar. \LaTeX{} es el estándar para documentos académicos, científicos y profesionales.

Este documento fue creado para probar tu setup de \LaTeX{} en VS Code. Si puedes ver el PDF generado, \textbf{¡todo funciona correctamente!}

% ── SECCIÓN 2 ──────────────────────────────
\section{Formato de Texto}

Puedes usar diferentes estilos de texto:

\begin{itemize}
    \item \textbf{Texto en negrita}
    \item \textit{Texto en cursiva}
    \item \underline{Texto subrayado}
    \item \texttt{Texto monoespaciado} (para código)
    \item \textsc{Texto en versalitas}
    \item Texto en \textcolor{red}{rojo}, \textcolor{blue}{azul} y \textcolor{green!50!black}{verde}
\end{itemize}

% ── SECCIÓN 3 ──────────────────────────────
\section{Listas}

\subsection{Lista no ordenada}
\begin{itemize}[label=\textbullet]
    \item Primer elemento
    \item Segundo elemento
    \item Tercer elemento con sub-items:
    \begin{itemize}
        \item Sub-item A
        \item Sub-item B
    \end{itemize}
\end{itemize}

\subsection{Lista ordenada}
\begin{enumerate}
    \item Paso uno: Instalar MiKTeX
    \item Paso dos: Instalar LaTeX Workshop
    \item Paso tres: Escribir tu documento
    \item Paso cuatro: Compilar y disfrutar
\end{enumerate}

% ── SECCIÓN 4 ──────────────────────────────
\section{Matemáticas}

\LaTeX{} es famoso por sus ecuaciones. Aquí algunos ejemplos:

\subsection{Ecuaciones en línea}
La famosa ecuación de Einstein: $E = mc^2$. La fórmula cuadrática es $x = \frac{-b \pm \sqrt{b^2 - 4ac}}{2a}$.

\subsection{Ecuaciones destacadas}
La integral de Gauss:
\begin{equation}
    \int_{-\infty}^{\infty} e^{-x^2} \, dx = \sqrt{\pi}
\end{equation}

La serie de Taylor:
\begin{equation}
    f(x) = \sum_{n=0}^{\infty} \frac{f^{(n)}(a)}{n!} (x - a)^n
\end{equation}

Un sistema de ecuaciones:
\begin{align}
    3x + 2y &= 7 \\
    x - y &= 1
\end{align}

% ── SECCIÓN 5 ──────────────────────────────
\section{Tablas}

\begin{table}[h]
    \centering
    \caption{Comparación de formatos de documento}
    \label{tab:formatos}
    \begin{tabular}{@{}lccc@{}}
        \toprule
        \textbf{Característica} & \textbf{\LaTeX} & \textbf{Word} & \textbf{Google Docs} \\
        \midrule
        Ecuaciones          & Excelente  & Bueno    & Básico  \\
        Tipografía          & Excelente  & Bueno    & Bueno   \\
        Control de versiones & Excelente & Limitado & Bueno   \\
        Colaboración        & Medio      & Bueno    & Excelente \\
        Curva de aprendizaje & Alta      & Baja     & Baja    \\
        \bottomrule
    \end{tabular}
\end{table}

% ── SECCIÓN 6 ──────────────────────────────
\section{Código Fuente}

Puedes incluir código con resaltado de sintaxis:

\begin{lstlisting}[language=Python, caption={Ejemplo en Python}]
def fibonacci(n):
    """Calcula el n-esimo numero de Fibonacci."""
    if n <= 1:
        return n
    a, b = 0, 1
    for _ in range(2, n + 1):
        a, b = b, a + b
    return b

# Imprimir los primeros 10 numeros
for i in range(10):
    print(f"F({i}) = {fibonacci(i)}")
\end{lstlisting}

% ── SECCIÓN 7 ──────────────────────────────
\section{Referencias y Enlaces}

Puedes hacer referencia a la Tabla~\ref{tab:formatos} fácilmente. También puedes agregar enlaces como \href{https://www.latex-project.org/}{la página oficial de \LaTeX}.

Para citas bibliográficas, puedes usar BibTeX (revisa la documentación para más detalles).

% ── CONCLUSIÓN ─────────────────────────────
\section{Conclusión}

¡Tu setup de \LaTeX{} está funcionando! Ahora puedes:

\begin{enumerate}
    \item Editar este archivo o crear nuevos \texttt{.tex}
    \item Compilar a PDF con \textbf{Ctrl+S} (se compila automáticamente)
    \item Exportar a Word con el task \textbf{``Exportar a Word''}
    \item Usar el script \texttt{build.ps1} para automatizar todo
\end{enumerate}

\vfill
\begin{center}
    \rule{0.5\textwidth}{0.4pt}\\[0.5em]
    \textit{Documento generado con \LaTeX{} desde VS Code}
\end{center}

\end{document}
